\section{Evaluierung Vincent mit Mitbewohner}

\subsection{Durchf\"uhrung}

\begin{itemize}
    \item Die Testperson war mein Mitbewohner und nicht Teil des Projektteams.
    \item Zu Beginn wurden alle vier Low-Fidelity-Prototypen vorgestellt: \textit{Data Feed}, \textit{Karten-App}, \textit{Story Explorer} und \textit{Wien in Zahlen}.
    \item Zu jedem Prototyp wurden die Grundidee, die wichtigsten Screens und der geplante Nutzungsablauf kurz erkl\"art.
    \item Anschlie{\ss}end wurde die Testperson nach ihrem allgemeinen Eindruck, ihren Pr\"aferenzen und m\"oglichen Verbesserungsw\"unschen gefragt.
    \item Der Fokus lag auf Verst\"andlichkeit, Zugang zu Open Data, hilfreichen Funktionen und unklaren Elementen.
\end{itemize}

\subsection{Ergebnisse}

\begin{itemize}
    \item Positiv wurde bewertet, dass alle vier Prototypen unterschiedliche Zug\"ange zu Open Data anbieten.
    \item Der \textit{Data Feed} wurde als vertraut wahrgenommen, da er an Social-Media-Feeds erinnert.
    \item Die \textit{Karten-App} wurde besonders f\"ur ortsbezogene Informationen als sinnvoll eingesch\"atzt.
    \item Beim \textit{Story Explorer} wurde der gef\"uhrte Ablauf von Thema zu Datensatz zu Story positiv hervorgehoben.
    \item \textit{Wien in Zahlen} wurde positiv bewertet, da Themenauswahl, personalisierter Feed und Fact Sheets einen klaren und alltagsnahen Zugang bieten.
    \item Besonders hilfreich wirkten bei \textit{Wien in Zahlen} das ``Grafik verstehen''-Element und der Abschnitt ``Was bedeutet das?''.
    \item Als bevorzugte Ans\"atze wurden die \textit{Karten-App} f\"ur konkrete Informationen in der Umgebung und \textit{Wien in Zahlen} f\"ur regelm\"a{\ss}iges, einfaches Informieren genannt.
    \item Als Verbesserungsw\"unsche wurden klarere Hinweise zur Aktualit\"at der Daten, sichtbare Quellenangaben und mehr Filterm\"oglichkeiten genannt.
\end{itemize}